\documentclass[10pt,a4paper]{amsart}
\usepackage[margin=19mm]{geometry}
\usepackage{amsmath,amssymb,mathtools}
\usepackage[scr=rsfs,cal=euler]{mathalfa}
\usepackage{xfrac}
\usepackage[unicode,hidelinks,bookmarks=false]{hyperref}
\newcommand{\ie}{i.e.\ }
\newcommand{\cf}{cf.\ }
\newcommand{\resp}{respectively\ }
\newcommand{\Subsection}{\subsection}
\providecommand{\nobreakdash}{\nobreak\hbox{-}}
\setlength{\emergencystretch}{2em}
\multlinegap0pt
\usepackage{amssymb}
\usepackage{dsfont}
\usepackage[normalem]{ulem}
\usepackage[shortlabels]{enumitem}
\usepackage{xcolor}
\usepackage{tikz}
\newtheorem{lemma}{Lemma}
\newtheorem{definition}{Definition}
\newtheorem{remark}{Remark}
\newtheorem{theorem}{Theorem}
\title{Thompson's Group $V$ and Virtual Link Theory}
\author{Micah Chrisman, Louisa Liles, Melody Molander}
\date{Source: arXiv:2607.28406v1 (CC BY 4.0)}
\begin{document}
\maketitle

\noindent\emph{Source and licence.} Thompson's Group $V$ and Virtual Link Theory. Version arXiv:2607.28406v1 (CC BY 4.0), distributed by arXiv under the Creative Commons Attribution 4.0 International licence, http://creativecommons.org/licenses/by/4.0/, as recorded in the arXiv licence metadata for this exact version and shown on the abstract page https://arxiv.org/abs/2607.28406v1. Full licence notice: \texttt{licenses/arxiv-2607.28406-CC-BY-4.0.txt}.\par\smallskip

\noindent\textit{Editorial scope.} Counted unit: the article's theorem function\_of\_positive\_type\_oriented (source0.tex lines 1162--1165). The article's complete original proof of it is delivered with it (source0.tex lines 1167--1185, one \texttt{proof} environment of 19 lines). No mathematics is written, repaired or re-derived: the statement and the proof are the article's own lines, in the article's own notation, and no retained line is substituted or re-flowed.  Retained uncounted as the source-local context the proof needs: lines 1061--1067; lines 1073--1079; lines 1081--1083; lines 1124--1130.  Nothing was omitted from any retained span, so the package carries no omission record.  No retained line contains a citation, so the wrapper carries no bibliography and no bibliography entry.  No retained line contains a figure environment, an image-inclusion command or drawing code, so no image is delivered and the package declares no asset dependency.  Editorial formatting only, in addition: an AMS \texttt{amsart} preamble that restates the article's own declarations and macros used by the retained lines, namely the environments \texttt{lemma}, \texttt{definition}, \texttt{remark}, \texttt{theorem} on separate counters, together with the standard packages those lines need.  The retained environments print in this document as Lemma 1, Definition 1, Remark 1, Theorem 1 in order of appearance, whereas the article prints the same items with the numbering of its own full document.  The complete native source of the article as distributed in the arXiv e-print for this exact version is bundled unchanged as \texttt{sources/206414/source0.tex}.
\begin{lemma}\label{lem: numberofopqcoloringswithcarets}
    Consider a triple $(T_+, T_-, \sigma)$ of two trees with $n$ leaves and an element $\sigma \in S_n$ representing an element $v \in \vec{V}$. Let $(S_+, S_-, \tau)$ be the new triple of trees with $n+1$ leaves and $\tau \in S_{n+1}$ obtained by introducing a canceling caret. Then
    \begin{align*}
        \text{Col}(\mathscr{L}_{V}(S_+,S_-,\tau), (Q, *,\alpha))=|Q|\text{Col}(\mathscr{L}_{V}(T_+,T_-,\sigma), (Q,*,\alpha)),
    \end{align*}
    for any finite operator quandle $(Q,*,\alpha)$.
\end{lemma}

\begin{definition}\label{def: oq-positive-type}
    Fix a finite operator quandle $(Q,*,\alpha)$. The $(Q,*,\alpha)$-\textit{coloring function} on $\vec{V}$, $\text{Col}(\cdot, (Q,*,\alpha)):\vec{V}\to \mathbb{R}_{\geq 0}$, is defined as
    \begin{align*}
        \text{Col}(v, (Q,*,\alpha)) := |Q|^{-n}\text{Col}(\mathscr{L}_{V}(T_+, T_-,\sigma), (Q,*,\alpha)),
    \end{align*}
    where $v \in \vec{V}$ is represented by the triple $(T_+, T_-,\sigma)$, $T_{\pm}$ have $n$ leaves, and $\sigma\in S_n$. 
\end{definition}

\begin{remark}\label{rmk: well-defined-positive-function}
    By Lemma \ref{lem: numberofopqcoloringswithcarets}, the addition of the normalization factor of $|Q|^{-n}$ in Definition \ref{def: oq-positive-type} makes the $(Q,*,\alpha)$-coloring function well-defined up to canceling carets. 
\end{remark}

\begin{remark}\label{rem: opq_colorings_is_Z}
    For $v \in \vec{V}$, the corresponding oriented virtual link diagram $\mathscr{L}_V(T_+,T_-,\sigma)$ satisfies: 
    \begin{align*}
        \text{Col}(\mathscr{L}_{V}(T_+,T_-,\sigma),(Q,*,\alpha))=Z_\Gamma(Q,*,\alpha),
    \end{align*}
    where $\Gamma$ is the 4-valent graph obtained from $\mathscr{L}_{V}(T_+,T_-,\sigma)$. 
\end{remark}

\begin{theorem}\label{thm: function_of_positive_type_oriented}
    The function $\text{Col}(\cdot,(Q,*,
\alpha)):\vec{V}\to \mathbb{R}_{\geq 0}$, where $(Q,*,\alpha)$ is any finite operator quandle is positive definite on $\vec{V}$.
\end{theorem}

\begin{proof}
    Fix a finite operator quandle $(Q,*,\alpha)$. Choose $r\in \mathbb{N}$ and $v_1,...,v_r \in \vec{V}$. We need to show that $(\text{Col}(v_iv_j^{-1}, (Q,*,\alpha)))_{1\leq i,j \leq r}$ is positive semidefinite. For any $1\leq i,j\leq r$, define $v_i:=\mathscr{L}_{V}(T_+^i,T_-^i, \sigma_i)$ and $v_j:=\mathscr{L}_{V}(T_+^j, T_-^j, \sigma_j)$ to be the virtual link diagrams representing $v_i$ and $v_j$ respectively where $T_+^i, T_-^i, T_+^j, T_-^j$ are four trees with $n$ leaves and $\sigma_i, \sigma_j \in S_n$, where $n$ is large enough to satisfy that for all $1\leq i,j \leq r$, $T_-^i=T_-^j$. Then $v_iv_j^{-1}=\mathscr{L}_{V}(T_+^i, T_+^j, \sigma_j^{-1} \sigma_i)$. By Remark \ref{rmk: well-defined-positive-function}, it suffices to show $\left(\text{Col}(\mathscr{L}_{V}(T_+^i, T_+^j, \sigma_j^{-1}\sigma_i), (Q,*,\alpha))\right)_{1\leq i,j\leq r}$ is positive semidefinite. 

    Let $\Gamma_{i,j}$ be the 4-valent graph created from the virtual link diagram $\mathscr{L}_{V}(T_+^i, T_+^j, \sigma_j^{-i}\sigma_i)$. We can decompose any function $\tau:E(\Gamma_{i,j})\to Q$ as $\tau=(\tau_0, \tau_+, \tau_-)$ where $\tau_0$ is the function on the middle $2n$ edges, and $\tau_\pm$ are the functions on the edges of the virtual semi-links (except the middle). Using Remark \ref{rem: opq_colorings_is_Z}, we can rewrite each entry of our matrix as:

    \begin{align*}
        &\text{Col}(\mathscr{L}_{V}(T_+^i,T_+^j, \sigma_j^{-1}\sigma_i, (Q,*,\alpha))=Z_\Gamma(Q,*,\alpha)\\
        &= \sum_{\tau_0} \Biggl(\sum_{\tau_+}\Biggl(\prod_{\substack{c \in \mathcal{C}_+^i\\ c\in \mathcal{C}_-^i} }\omega_{c} (a_{\text{in}}^\text{over}, b_{\text{in}}^\text{under}, a_{\text{out}}^\text{over},b_{\text{out}}^\text{under})\prod_{c_\text{o}\in \mathcal{C}_\text{o}^i}\omega_{c_\text{o}}(a_\text{in}, b_\text{in}, a_\text{out}, b_\text{out})\Biggr)\\
        &\times \sum_{\tau_-}\Biggl(\prod_{\substack{c \in \mathcal{C}_+^j\\ c\in \mathcal{C}_-^j} }\omega_{c} (a_{\text{in}}^\text{over}, b_{\text{in}}^\text{under}, a_{\text{out}}^\text{over},b_{\text{out}}^\text{under})\prod_{c_\text{o}\in \mathcal{C}_\text{o}^j}\omega_{c_\text{o}}(a_\text{in}, b_\text{in}, a_\text{out}, b_\text{out})\Biggr)\Biggr).
    \end{align*}
    The first sum, over $\tau_0$, is over the $|Q|^{2n}$ functions coloring the $2n$ strands in the middle of the two virtual semi-links. The sums over $\tau_+, \tau_-$ are over the functions coloring the strands in the respective upper and lower virtual semi-link. The sets $\mathcal{C}_\pm^i, \mathcal{C}_o^i$ now only contain the crossings from the $T_+^i$ tree (likewise for $j$). 

    Since $\tau_0$ is a function on each of the $2n$ strands in the middle, we can write $\tau_0=(\tau_1,...,\tau_{2n})\in |Q|^{2n}$. Recall that $\ell^2(|Q|^{2n})=\{f:|Q|^{2n} \to \mathbb{C}:\sum_{\tau_0\in |Q|^{2n}} |f(\tau_0)|^2<\infty\}$ is a Hilbert space with inner product $\langle f, g\rangle =\sum_{\tau_0 \in |Q|^{2n}} f \overline{g}$. This Hilbert space is isomorphic to $\mathbb{C}^{|Q|^{2n}}$. We claim that $f_i:|Q|^{2n}\to \mathbb{C}$ defined by:
    \begin{align*}
        \tau_0 \mapsto \sum_{\tau_+}\left(\prod_{\substack{c \in \mathcal{C}_+^i\\ c\in \mathcal{C}_-^i} }\omega_{c} (a_{\text{in}}^\text{over}, b_{\text{in}}^\text{under}, a_{\text{out}}^\text{over},b_{\text{out}}^\text{under})\prod_{c_\text{o}\in \mathcal{C}_\text{o}^i}\omega_{c_\text{o}}(a_\text{in}, b_\text{in}, a_\text{out}, b_\text{out})\right),
    \end{align*}
    is in $\ell^2(|Q|^{2n})$. Indeed, $\sum_{\tau_0\in |Q|^{2n}}|f_i(\tau_0)|^2=\sum_{\tau_0\in |Q|^{2n}}f_i(\tau_0)^2$ is finite since the sum itself is finite. Further, for any $1\leq i,j \leq r$, $\langle f_i, f_j\rangle=Z_\Gamma(Q,*,\alpha)=\text{Col}(\mathscr{L}_{V}(T_+^i,T_+^j, \sigma_j^{-1}\sigma_i, (Q,*,\alpha))$, as we wished. By taking an $r$-dimensional closed subspace containing $f_1,...,f_r$ in $\ell^2(|Q|^{2n})$, we conclude that the matrix $\left(\text{Col}(\mathscr{L}_{V}(T_+^i, T_+^j, \sigma_j^{-1}\sigma_i), (Q,*,\alpha))\right)_{1\leq i,j \leq r}$ is positive semidefinite and further that the function $\text{Col}(\cdot,(Q,*,
\alpha))$ is positive definite on $\vec{V}$.
\end{proof}

\end{document}
